%%%%%%%%%%%%%%%%%%%%%%%%%%%%%%%%%%%%%%%%%%%%%%%%%%%%%%%%%%%%%%
\chapter{Conclusion}
\label{chp:conclusion}
%%%%%%%%%%%%%%%%%%%%%%%%%%%%%%%%%%%%%%%%%%%%%%%%%%%%%%%%%%%%%%
\begin{hints}
\item One page summary of what was done: data, model construction, layer-wise validation, scenario.
\end{hints}

\section{Answers to the Research Questions}
\label{sec:concl_rq}
\begin{hints}
\item One paragraph per RQ, in the pattern: the question in one sentence, the direct answer, the key evidence (one or two numbers), the qualification. Same RQ wording as in Chapter~\ref{chp:introduction}.
\item RQ1: take the short answer of Section~\ref{sec:res_rq1}. Part~1: IM 0.0 restricted exactly to rules 3 and 5 as a state machine, plus the stop guard; the only sound net that accepts every observed sequence with precision clearly above a flower-like net (0.838, repaired 0.888); preserving rare behaviour is not overfitting (13 of 16 new test sequences accepted). Part~2: only partly; precision did not anticipate the simulation (same language and precision, but 15.1\,\% vs.\ 0\,\% visits without a stop), coverage did (IMf 0.3: lowest coverage, only clearly worse structure); the state-machine encoding decided.
\item RQ2: take the short answer of Section~\ref{sec:res_rq2}. As fitted, no (median 58 vs.\ 33 min, P90 141 vs.\ 75); with replayed trucks and blocks and four changes (exact count, capacity at start with three lanes, lane-free waiting target, load-dependent service) to a large extent: median 28, P90 62 min on the test period, in-sample CTD 2.7 min (within the variation between weeks); approach rises with the block load as observed. Causes in the order of the truck: block choice (self-count in the queue feature, amplified by the too low concurrency count) for the E3 error; capacity at enabling and single-server waiting for the missing waiting under load; no load in the service time. Qualification: underestimates the slowdown at the busiest blocks and hours; the load-response criterion is missed narrowly, so E5 is not admissible by the pre-registered rule; part of the test error is the slower service of the test period.
\item RQ3: take the short answer of the E5 results. What the simulation can and cannot say about load-aware block allocation for deliveries: direction and mechanism only if robust in all four robustness runs; effect sizes as model values under the conditions of E4; no recommendation for the terminal. Frame it as a demonstration of what-if use of process mining based simulation under physical constraints, with its conditions of validity.
\end{hints}

\section{Contributions}
\label{sec:concl_contributions}
\begin{hints}
\item Restate the contributions from Chapter~\ref{chp:introduction}, now with the evidence that supports each.
\end{hints}

\section{Limitations}
\label{sec:concl_limitations}
\begin{hints}
\item The three or four most important limitations, short (details in Chapter~\ref{chp:discussion}).
\end{hints}

%%%%%%%%%%%%%%%%%%%%%%%%%%%%%%%%%%%%%%%%%%%%%%%%%%%%%%%%%%%%%%
\section{Future Work}
\label{sec:future_work}
%%%%%%%%%%%%%%%%%%%%%%%%%%%%%%%%%%%%%%%%%%%%%%%%%%%%%%%%%%%%%%

Future work should retain the reproducible discovery-and-validation protocol
while extending the mechanisms and evidence that currently limit physical and
counterfactual interpretation:

\begin{itemize}
    \item \textbf{Investigate an object-centric process representation.}
    Future work could examine whether an object-centric event log is better
    suited to container-terminal operations than the current truck-visit case
    representation. In the present model, container characteristics are
    aggregated to case level, so information such as container type, loading
    status, reefer status, or hazardous-goods status is combined across all
    containers handled during one truck visit. An object-centric approach could
    keep each container as a separate object with its own attributes while
    linking it to the truck visit, handling activity, yard block, and resource.
    This would reduce information loss and preserve container-specific information throughout the process
    and could also support multi-container and multi-block visits without
    forcing one aggregated context value onto the complete case.

    \item \textbf{Add an explicit spatial and movement representation.}
    Future work could separate physical movement from the combined
    pre-service delay by using route, position, or equipment observations.
    Rather than representing every movement as an additional process activity,
    a richer simulation formalism could include location directly in the model
    state, for example through timed or coloured Petri-net tokens and
    location-dependent transition conditions. This would allow truck or
    container movement between yard areas to consume time and respect spatial
    constraints without treating travel as a recorded handling activity.
    Movement times would then have to be estimated from spatial observations
    and the remaining access and waiting-time models refitted, so that travel
    time is not counted twice.

   \item \textbf{Represent physical assignment and queueing explicitly.}
    Future models could bind trucks and containers to physically feasible
    yard locations and resources instead of treating resource labels as fully
    interchangeable. Equipment availability and assignment should be represented
    as dynamic state, and trucks waiting for handling should enter explicit
    simulated queues rather than having these effects absorbed into a fitted
    pre-service-delay distribution. An object-centric representation could
    preserve the relations between trucks, containers, yard blocks, and
    equipment, while a corresponding simulation model would still need explicit
    movement, resource-allocation, and queueing rules
    \cite{VanDerAalst2023OCPM,Knopp2023ObjectCentricSimulation}.

    
    
\end{itemize}
